\documentclass[11pt,a4paper]{article}

% Packages
\usepackage[utf8]{inputenc}
\usepackage[T1]{fontenc}
\usepackage{lmodern}
\usepackage[margin=1in]{geometry}
\usepackage{amsmath,amssymb,amsthm}
\usepackage{graphicx}
\usepackage{hyperref}
\usepackage{xcolor}
\usepackage{listings}
\usepackage{algorithm}
\usepackage{algpseudocode}
\usepackage{booktabs}
\usepackage{tikz}
\usetikzlibrary{shapes,arrows,positioning}

% Colors
\definecolor{tsnblue}{RGB}{59,130,246}
\definecolor{codegray}{RGB}{245,245,245}
\definecolor{quantumgreen}{RGB}{34,197,94}

% Hyperref setup
\hypersetup{
    colorlinks=true,
    linkcolor=tsnblue,
    urlcolor=tsnblue,
    citecolor=tsnblue
}

% Code listing style
\lstset{
    basicstyle=\ttfamily\small,
    backgroundcolor=\color{codegray},
    frame=single,
    breaklines=true
}

% Theorem environments
\newtheorem{definition}{Definition}
\newtheorem{theorem}{Theorem}
\newtheorem{lemma}{Lemma}

% Title
\title{
    \vspace{-1cm}
    \textbf{TSN: A Fully Quantum-Resistant Privacy Cryptocurrency}\\[0.5em]
    \large Technical Whitepaper v2.0
}
\author{
    Trust Stack Network\\
    \texttt{https://tsnchain.com}
}
\date{February 2026}

\begin{document}

\maketitle

\begin{abstract}
We present TSN v2, a cryptocurrency providing strong privacy guarantees that remain secure against both classical and quantum adversaries. Unlike previous versions that used quantum-vulnerable pairing-based proofs for privacy, TSN v2 achieves \textbf{full post-quantum security} by combining NIST-standardized signatures (ML-DSA-65, FIPS 204) with hash-based STARK proofs (Plonky3 over the Goldilocks field). This eliminates all reliance on elliptic curve discrete logarithm assumptions, providing complete protection against Shor's algorithm. The protocol maintains a note-based UTXO model similar to Zcash's Sapling but replaces all quantum-vulnerable components with lattice-based and hash-based alternatives. This paper describes the cryptographic foundations, V2 transaction model, migration path from V1, and security properties of the fully quantum-resistant TSN protocol.
\end{abstract}

\tableofcontents
\newpage

%==============================================================================
\section{Introduction}
%==============================================================================

\subsection{The Quantum Threat}

The development of practical quantum computers poses an existential threat to current cryptographic systems. Shor's algorithm can solve the discrete logarithm problem in polynomial time, breaking:

\begin{itemize}
    \item \textbf{ECDSA/EdDSA signatures}: Used by Bitcoin, Ethereum, and most cryptocurrencies
    \item \textbf{Pairing-based cryptography}: Used by Groth16 zk-SNARKs (Zcash, Tornado Cash)
    \item \textbf{Pedersen commitments}: Used for hiding transaction amounts
\end{itemize}

While TSN v1 addressed signature security with ML-DSA-65, the privacy layer remained vulnerable due to its use of BN254-based Groth16 proofs and Pedersen commitments.

\subsection{TSN v2: Full Quantum Resistance}

TSN v2 eliminates all quantum-vulnerable components:

\begin{table}[h]
\centering
\begin{tabular}{llll}
\toprule
Component & V1 & V2 & Quantum Safe \\
\midrule
Signatures & ML-DSA-65 & ML-DSA-65 & \textcolor{quantumgreen}{\checkmark} \\
ZK Proofs & Groth16/BN254 & Plonky3 STARKs & \textcolor{quantumgreen}{\checkmark} \\
Commitments & Pedersen/BN254 & Poseidon/Goldilocks & \textcolor{quantumgreen}{\checkmark} \\
Merkle Tree & Poseidon/BN254 & Poseidon/Goldilocks & \textcolor{quantumgreen}{\checkmark} \\
Binding Sig & Schnorr/BN254 & Included in STARK & \textcolor{quantumgreen}{\checkmark} \\
\bottomrule
\end{tabular}
\caption{Comparison of V1 and V2 cryptographic components}
\end{table}

%==============================================================================
\section{Threat Model and Security Goals}
%==============================================================================

\subsection{Adversary Capabilities}

We consider adversaries with the following capabilities:

\begin{itemize}
    \item \textbf{Classical Adversary}: Polynomial-time bounded classical computer with access to all public blockchain data.
    \item \textbf{Quantum Adversary}: Access to a cryptographically-relevant quantum computer capable of running Shor's algorithm (breaks discrete log) and Grover's algorithm (quadratic speedup for search).
    \item \textbf{Network Adversary}: Ability to observe, delay, or reorder network messages.
\end{itemize}

\subsection{Security Goals}

\begin{definition}[Post-Quantum Transaction Integrity]
No adversary, even with quantum computing capabilities, can create a valid transaction that spends funds they do not own.
\end{definition}

\begin{definition}[Post-Quantum Balance Privacy]
No adversary, even with quantum computing capabilities, can determine the balance of any address without knowledge of the corresponding viewing key.
\end{definition}

\begin{definition}[Post-Quantum Transaction Privacy]
No adversary, even with quantum computing capabilities, can link the sender and receiver of a transaction, nor determine the transaction amount.
\end{definition}

%==============================================================================
\section{Cryptographic Primitives}
%==============================================================================

\subsection{Post-Quantum Digital Signatures: ML-DSA-65}

TSN uses ML-DSA-65 (Module-Lattice Digital Signature Algorithm), standardized as FIPS 204 by NIST.

\begin{table}[h]
\centering
\begin{tabular}{lrl}
\toprule
Parameter & Value & Description \\
\midrule
Security Level & 3 (128-bit) & NIST security category \\
Public Key Size & 1,952 bytes & Verification key \\
Secret Key Size & 4,032 bytes & Signing key \\
Signature Size & 3,309 bytes & Digital signature \\
\bottomrule
\end{tabular}
\caption{ML-DSA-65 parameters}
\end{table}

The security of ML-DSA relies on the hardness of Module-LWE and Module-SIS problems, which are believed hard for both classical and quantum computers.

\subsection{Post-Quantum Zero-Knowledge Proofs: Plonky3 STARKs}

TSN v2 replaces Groth16 with Plonky3, a STARK-based proving system that relies only on hash function security.

\begin{definition}[STARK]
A Scalable Transparent ARgument of Knowledge is a proof system where:
\begin{itemize}
    \item \textbf{Transparency}: No trusted setup required
    \item \textbf{Post-Quantum Security}: Security relies on collision-resistant hash functions
    \item \textbf{Scalability}: Prover time quasi-linear, verifier time polylogarithmic
\end{itemize}
\end{definition}

\subsubsection{Plonky3 Parameters}

\begin{table}[h]
\centering
\begin{tabular}{ll}
\toprule
Parameter & Value \\
\midrule
Field & Goldilocks ($p = 2^{64} - 2^{32} + 1$) \\
Hash Function & Poseidon (width-12, rate-8) \\
FRI Parameters & Rate 1/8, 4 query rounds \\
Security Level & $\geq 100$ bits post-quantum \\
Proof Size & $\approx 50$ KB \\
\bottomrule
\end{tabular}
\caption{Plonky3 configuration in TSN v2}
\end{table}

\subsubsection{Why Plonky3?}

\begin{enumerate}
    \item \textbf{Hash-based security}: No reliance on discrete log or pairings
    \item \textbf{No trusted setup}: Unlike Groth16, no toxic waste
    \item \textbf{Fast proving}: Optimized for client-side browser execution
    \item \textbf{Recursive composition}: Enables future protocol upgrades
\end{enumerate}

\subsection{Hash Functions}

\subsubsection{Poseidon over Goldilocks}

All in-circuit hashing uses Poseidon optimized for the Goldilocks field:

\begin{equation}
H_{\text{Poseidon}}: \mathbb{F}_p^n \rightarrow \mathbb{F}_p^4
\end{equation}

where $p = 2^{64} - 2^{32} + 1$ (Goldilocks prime).

The hash outputs 4 field elements (256 bits total), matching Plonky3's internal representation.

\begin{table}[h]
\centering
\begin{tabular}{ll}
\toprule
Domain & Separator \\
\midrule
Note Commitment & 1 \\
Value Commitment & 2 \\
Nullifier & 3 \\
Merkle Empty Leaf & 4 \\
Merkle Node & 5 \\
\bottomrule
\end{tabular}
\caption{V2 Poseidon domain separators}
\end{table}

\subsubsection{Post-Quantum Security of Hash Functions}

Grover's algorithm provides a quadratic speedup for brute-force search, reducing effective security:

\begin{itemize}
    \item 256-bit hash $\rightarrow$ 128-bit post-quantum security
    \item Goldilocks field elements: 4 $\times$ 64 bits = 256 bits
\end{itemize}

This provides adequate post-quantum security margins.

%==============================================================================
\section{V2 Transaction Model}
%==============================================================================

\subsection{Notes}

The fundamental unit of value remains a \textbf{note}:

\begin{definition}[V2 Note]
A note $\mathcal{N}$ is a tuple $(v, pk_h, r)$ where:
\begin{itemize}
    \item $v \in \{0, 1, \ldots, 2^{64}-1\}$ is the value
    \item $pk_h \in \mathbb{F}_p^4$ is the recipient's public key hash (4 Goldilocks elements)
    \item $r \in \mathbb{F}_p^4$ is random blinding factor (4 Goldilocks elements)
\end{itemize}
\end{definition}

\subsection{V2 Commitments}

Unlike V1's Pedersen commitments (which hide values using elliptic curve discrete log), V2 uses hash-based commitments:

\begin{definition}[V2 Note Commitment]
The commitment to a note $\mathcal{N} = (v, pk_h, r)$ is:
\begin{equation}
\text{cm} = H_{\text{Poseidon}}(1, v, pk_h[0..3], r[0..3])
\end{equation}
Output: 4 Goldilocks field elements (32 bytes serialized).
\end{definition}

\textbf{Security}: The commitment is:
\begin{itemize}
    \item \textbf{Hiding}: Without $r$, the value cannot be determined (hash preimage resistance)
    \item \textbf{Binding}: Cannot find two different notes with the same commitment (collision resistance)
\end{itemize}

Both properties hold against quantum adversaries.

\subsection{V2 Nullifiers}

\begin{definition}[V2 Nullifier]
The nullifier for a note with commitment $\text{cm}$ at position $\rho$, using nullifier key $nk$:
\begin{equation}
\text{nf} = H_{\text{Poseidon}}(3, nk[0..3], \text{cm}[0..3], \rho)
\end{equation}
\end{definition}

Note: V2 nullifiers include the full commitment (not just position), providing stronger unlinkability.

\subsection{V2 Commitment Tree}

A Merkle tree of depth 32 using Poseidon/Goldilocks:

\begin{itemize}
    \item \textbf{Leaves}: Note commitments (4 field elements each)
    \item \textbf{Internal nodes}: $H_{\text{Poseidon}}(5, \text{left}[0..3], \text{right}[0..3])$
    \item \textbf{Empty leaf}: $H_{\text{Poseidon}}(4, 0, 0, 0, 0)$
\end{itemize}

\subsection{V2 Transaction Structure}

A V2 shielded transaction consolidates all proofs into a single STARK:

\begin{lstlisting}
ShieldedTransactionV2 {
    version: 2,
    spends: [{
        anchor: [u8; 32],      // Merkle root
        nullifier: [u8; 32],   // Spend nullifier
        signature: [u8; 3309], // ML-DSA-65 signature
        public_key: [u8; 1952] // ML-DSA-65 public key
    }],
    outputs: [{
        note_commitment: [u8; 32],
        encrypted_note: { ciphertext, ephemeral_pk }
    }],
    fee: u64,
    transaction_proof: {
        proof_bytes: Vec<u8>,  // Plonky3 STARK (~50KB)
        public_inputs: {
            merkle_roots: [[u8; 32]],
            nullifiers: [[u8; 32]],
            note_commitments: [[u8; 32]],
            fee: u64
        }
    }
}
\end{lstlisting}

\subsection{Combined STARK Proof}

The V2 STARK proves all of the following in a single proof:

\begin{enumerate}
    \item \textbf{For each spend}:
    \begin{itemize}
        \item Knowledge of note preimage $(v, pk_h, r)$
        \item Note commitment exists in tree (Merkle path verification)
        \item Nullifier correctly derived from $(nk, \text{cm}, \rho)$
    \end{itemize}

    \item \textbf{For each output}:
    \begin{itemize}
        \item Note commitment correctly computed from $(v, pk_h, r)$
    \end{itemize}

    \item \textbf{Balance constraint}:
    \begin{equation}
    \sum_i v_{\text{spend},i} = \sum_j v_{\text{output},j} + \text{fee}
    \end{equation}
\end{enumerate}

This eliminates the need for separate binding signatures---the balance is enforced inside the STARK.

\subsection{Dynamic Circuit Building}

Unlike V1's fixed circuits, V2 builds circuits dynamically:

\begin{itemize}
    \item Circuits are parameterized by $(n_{\text{spends}}, n_{\text{outputs}})$
    \item Maximum: 10 spends $\times$ 4 outputs
    \item Circuits are cached after first build for reuse
    \item ``Warmup'' pre-builds common shapes (1-5 spends $\times$ 1-2 outputs)
\end{itemize}

%==============================================================================
\section{V1 to V2 Migration}
%==============================================================================

\subsection{Migration Transaction}

Users with V1 notes can migrate to V2:

\begin{enumerate}
    \item \textbf{Spend V1 notes}: Using legacy Groth16 proofs
    \item \textbf{Create V2 outputs}: With Poseidon/Goldilocks commitments
    \item \textbf{Combined proof}: V1 spends + V2 outputs in single transaction
\end{enumerate}

\subsection{Migration Incentives}

\begin{itemize}
    \item V2 transactions have lower fees (no binding signature verification)
    \item V1 will be deprecated after sufficient migration period
    \item Quantum computers are advancing---migrate before they arrive
\end{itemize}

\subsection{Deprecation Timeline}

\begin{enumerate}
    \item \textbf{Phase 1} (Current): Both V1 and V2 supported, V2 default
    \item \textbf{Phase 2} (6 months): V1 creation disabled, only spending allowed
    \item \textbf{Phase 3} (12 months): V1 support removed entirely
\end{enumerate}

%==============================================================================
\section{Security Analysis}
%==============================================================================

\subsection{Full Post-Quantum Security}

\begin{theorem}[V2 Quantum Resistance]
TSN V2 transactions are secure against adversaries with access to cryptographically-relevant quantum computers.
\end{theorem}

\begin{proof}
V2 relies only on:
\begin{enumerate}
    \item \textbf{ML-DSA-65}: Based on Module-LWE/SIS, no known quantum speedup beyond Grover's $O(\sqrt{N})$
    \item \textbf{Poseidon hash}: Collision resistance reduced to $2^{128}$ operations with Grover
    \item \textbf{Plonky3 STARKs}: Security reduces to hash collision resistance
\end{enumerate}
No component relies on discrete log or factoring, which Shor's algorithm breaks.
\end{proof}

\subsection{Comparison with V1}

\begin{table}[h]
\centering
\begin{tabular}{lcc}
\toprule
Attack & V1 Vulnerable? & V2 Vulnerable? \\
\midrule
Forge signatures & No (ML-DSA) & No (ML-DSA) \\
Forge ZK proofs & Yes (BN254) & No (STARK) \\
Open commitments & Yes (Pedersen) & No (Poseidon) \\
Break binding sig & Yes (Schnorr/BN254) & N/A (in STARK) \\
Link transactions & Yes (if proofs forged) & No \\
\bottomrule
\end{tabular}
\caption{Quantum attack vectors}
\end{table}

\subsection{Double-Spend Prevention}

\begin{theorem}[Double-Spend Resistance]
A V2 note can only be spent once, even by a quantum adversary.
\end{theorem}

\begin{proof}
The nullifier $\text{nf} = H(nk, \text{cm}, \rho)$ is deterministic. The STARK proves correct nullifier derivation. Duplicate nullifiers are rejected by consensus. Hash collision would require $2^{128}$ quantum operations.
\end{proof}

%==============================================================================
\section{Performance}
%==============================================================================

\begin{table}[h]
\centering
\begin{tabular}{lrll}
\toprule
Operation & V1 Time & V2 Time & Notes \\
\midrule
Proof generation (1 spend) & 3s & 2s & Browser WASM \\
Proof generation (4 spends) & 12s & 5s & Browser WASM \\
Proof verification & 5ms & 50ms & Server-side \\
Proof size & 200 bytes & 50 KB & Per transaction \\
Circuit warmup (5$\times$2) & N/A & 30s & One-time \\
\bottomrule
\end{tabular}
\caption{Performance comparison}
\end{table}

The larger proof size is an acceptable trade-off for quantum resistance. Proof generation is actually \textit{faster} in V2 due to Plonky3's optimized WASM implementation.

%==============================================================================
\section{Implementation}
%==============================================================================

\subsection{Browser-Based WASM Prover}

The Plonky3 prover is compiled to WebAssembly for client-side execution:

\begin{lstlisting}
// TypeScript interface
const prover = new WasmProver();
await prover.prebuild_circuit(4, 2); // Warmup
const proof = prover.prove(witnessJson);
\end{lstlisting}

Key features:
\begin{itemize}
    \item Dynamic circuit building for any transaction shape
    \item Circuit caching for instant subsequent proofs
    \item Background warmup with progress tracking
    \item LocalStorage persistence of warmup state
\end{itemize}

\subsection{Server-Side Verification}

Proof verification uses native Rust Plonky3:

\begin{lstlisting}
let circuit = CircuitCache::get_or_build(n_spends, n_outputs);
circuit.verify(proof)?;
\end{lstlisting}

The circuit cache uses thread-safe \texttt{Arc<RwLock>} for concurrent access.

%==============================================================================
\section{Future Work}
%==============================================================================

\begin{enumerate}
    \item \textbf{Recursive Proofs}: Aggregate multiple transactions into single proof
    \item \textbf{Light Clients}: STARK-based SPV proofs for mobile
    \item \textbf{Proof Compression}: Reduce 50KB proofs using recursive SNARKs
    \item \textbf{Post-Quantum Key Exchange}: Replace ChaCha20-Poly1305 key derivation with Kyber
\end{enumerate}

%==============================================================================
\section{Conclusion}
%==============================================================================

TSN v2 achieves \textbf{complete post-quantum security} for private cryptocurrency transactions. By replacing all elliptic curve-based components with hash-based alternatives:

\begin{itemize}
    \item \textbf{Signatures}: ML-DSA-65 (lattice-based, NIST standardized)
    \item \textbf{ZK Proofs}: Plonky3 STARKs (hash-based, no trusted setup)
    \item \textbf{Commitments}: Poseidon/Goldilocks (hash-based)
\end{itemize}

Users can transact privately today with confidence that their transactions will remain private and secure even after the advent of large-scale quantum computers.

%==============================================================================
\section*{Acknowledgments}
%==============================================================================

We thank the developers of Plonky3 (Polygon Labs), the NIST PQC standardization team, and the Zcash protocol designers whose work inspired TSN's privacy model.

%==============================================================================
\begin{thebibliography}{99}

\bibitem{fips204}
National Institute of Standards and Technology.
\textit{FIPS 204: Module-Lattice-Based Digital Signature Standard}.
2024.

\bibitem{plonky2}
Polygon Labs Team.
\textit{Plonky3: A Toolkit for Succinct Proof Systems}.
2024.

\bibitem{stark}
Ben-Sasson, E., Bentov, I., Horesh, Y., Riabzev, M.
\textit{Scalable, Transparent, and Post-Quantum Secure Computational Integrity}.
IACR Cryptology ePrint Archive, 2018.

\bibitem{goldilocks}
Hamburg, M.
\textit{Ed448-Goldilocks, a New Elliptic Curve}.
IACR Cryptology ePrint Archive, 2015.

\bibitem{poseidon}
Grassi, L., Khovratovich, D., Rechberger, C., Roy, A., Schofnegger, M.
\textit{Poseidon: A New Hash Function for Zero-Knowledge Proof Systems}.
USENIX Security 2021.

\bibitem{zcash}
Hopwood, D., Bowe, S., Hornby, T., Wilcox, N.
\textit{Zcash Protocol Specification}.
2023.

\bibitem{shor}
Shor, P.
\textit{Algorithms for Quantum Computation: Discrete Logarithms and Factoring}.
FOCS 1994.

\bibitem{grover}
Grover, L.
\textit{A Fast Quantum Mechanical Algorithm for Database Search}.
STOC 1996.

\end{thebibliography}

\end{document}
